\section{Graph thresholds and the main result}\label{sec:introduction}

When does a random graph contain a prescribed graph? A necessary
condition is that each subgraph of the target have a reasonable expected
number of copies. Kahn and Kalai~\cite{KK07} conjectured
that this elementary obstruction determines the containment threshold
within a logarithmic factor, uniformly over the target graph. We prove
this conjecture, with the logarithm depending on the number of target
edges.

For finite simple graphs $G$ and $F$, write $N(G,F)$ for the number of
subgraphs of $G$ isomorphic to $F$. These are ordinary copies: additional
edges between the chosen vertices are allowed. Let $G(n,p)$ be the
random graph on $[n]=\{1,\ldots,n\}$ in which each edge is present
independently with probability $p$. For a graph $H$ with at least one
edge and at most $n$ vertices, define
\begin{align*}
 \pc(n,H)
 &=\inf\{p\in[0,1]:\PP(N(G(n,p),H)>0)\ge\tfrac12\},\\
 \pE(n,H)
 &=\inf\{p\in[0,1]:\EE N(G(n,p),F)\ge\tfrac12
                 \text{ for every subgraph }F\subseteq H\}.
\end{align*}
The first quantity is the containment threshold; the second is the
\emph{graph expectation threshold}. Both infima are attained, since the
relevant probabilities and expectations are continuous, there are only
finitely many constraints, and $p=1$ satisfies all of them. The
single-edge constraint gives $\pE(n,H)>0$. If $H$ appears with
probability at least one half, then so does every $F\subseteq H$, and
its expected count is at least one half. Thus $\pE(n,H)\le\pc(n,H)$.

\begin{theorem}\label{thm:main}
For every $n\ge2$ and every finite simple graph $H$ with $h=e(H)\ge1$
edges and at most $n$ vertices,
\begin{equation}\label{eq:stronger}
 \pc(n,H)\le
 \min\{1,\;2048\ee^{50}\pE(n,H)(1+\log_2 h)\}.
\end{equation}
In particular,
\[
 \pc(n,H)\le6144\ee^{50}\pE(n,H)\log_2 n.
\]
\end{theorem}

Here $\ee$ is the base of the natural logarithm. The constants are
explicit and universal; the graph $H$ may vary with $n$.
Kahn and Kalai use expected count one in their formulation. Replacing
one half by one changes the graph expectation threshold by a factor
between one and two: for each subgraph with $r\ge1$ edges its expected
count is proportional to $p^r$, so its constraint changes by
$2^{1/r}\le2$. This normalization therefore gives the same conjecture.

\subsection{History and the obstruction}

For a fixed target, threshold questions are part of the classical
random-graph theory of Erd\H{o}s and R\'enyi~\cite{ER60};
Bollob\'as~\cite{Bollobas81} treated threshold functions for general
fixed subgraphs. Here the target can grow with the host graph.
For example, let $H$ be a perfect matching on an even number $n$ of
vertices. Then $\pE(n,H)=\Theta(1/n)$, whereas
$\pc(n,H)=\Theta(\log n/n)$~\cite{ER66,KK07}.
The larger scale is needed to eliminate isolated vertices of the host:
a perfect matching must cover every vertex. Hamilton cycles exhibit
the same logarithmic gap~\cite{KK07}. These examples show why a
uniform comparison must allow a logarithm. Other spanning targets have
required different methods: Riordan~\cite{Riordan00}, for example,
proved that for every fixed $p>1/4$, $G(2^d,p)$ contains the
$d$-dimensional cube with probability tending to one.

To distinguish the two Kahn--Kalai conjectures, consider a nonempty
proper increasing family $\mathcal F\subseteq2^X$, where $X$ is finite.
A \emph{cover} is a family $\mathcal G\subseteq2^X$ such that every
$F\in\mathcal F$ contains some $S\in\mathcal G$. Its cost at density
$p$ is $\sum_{S\in\mathcal G}p^{|S|}$. The integral expectation
threshold $q(\mathcal F)$ is the largest $p$ for which a cover has
cost at most $1/2$. A fractional cover assigns weights $w_S\ge0$
to $S\subseteq X$, requires $\sum_{S\subseteq F}w_S\ge1$ for each
$F\in\mathcal F$, and has cost $\sum_S w_Sp^{|S|}$. Allowing these
covers defines $q_f(\mathcal F)$ with the same budget $1/2$.
Union bounds and expectations give
$q(\mathcal F)\le q_f(\mathcal F)\le p_{\mathrm c}(\mathcal F)$,
where $p_{\mathrm c}(\mathcal F)$ is the density at which an independent
random subset belongs to $\mathcal F$ with probability $1/2$.

Talagrand~\cite[Section~6]{Tal10} introduced the fractional formulation
and its dual description by spread probability measures.
Alweiss, Lovett, Wu, and Zhang~\cite{ALWZ21}, building on the regularity
viewpoint of Lovett, Solomon, and Zhang~\cite{LSZ19}, developed the
fragment method in their work on sunflowers. Frankston, Kahn,
Narayanan, and Park~\cite{FKNP21} proved
$p_{\mathrm c}(\mathcal F)=O(q_f(\mathcal F)\log\ell)$, where $\ell$
is the maximum of two and the largest size of a minimal member of
$\mathcal F$. Park and Pham~\cite{PP24} proved the corresponding
bound with $q$ in place of $q_f$, resolving the abstract Kahn--Kalai
conjecture by a minimum-fragment argument.
The integral and fractional thresholds are within a universal constant
factor at the common one-half budget
\cite[Theorem~1.1]{OpenAIThresholdEquivalence2026}.
This comparison is contextual; the proof below does not use it.

For the family $\mathcal F_H$ of edge sets containing a copy of $H$,
the copies of any one subgraph form a cover. Thus
$\pE(n,H)\le q(\mathcal F_H)$, and the abstract theorem alone does
not give the desired bound in terms of $\pE(n,H)$.
Mossel, Niles-Weed, Sun, and Zadik~\cite[Theorem~1]{MNSZ22} obtained
a single logarithm using the modified graph threshold
\[
 \widetilde p_{\mathrm E}(n,H)
 =\min\{p:\EE N(G(n,p),F)\ge N(H,F)/2
                  \text{ for every }F\subseteq H\}.
\]
The factor $N(H,F)$ records that containing $H$ supplies all its
copies of $F$. By symmetry this parameter equals
$q_f(\mathcal F_H)$~\cite[Equation~(1)]{DKP25}; comparing it with
$\pE(n,H)$ is the additional graph-specific problem.

Dubroff, Kahn, and Park~\cite[Theorems~1.2--1.5]{DKP25} proved
$q_f(\mathcal F_H)=O(\pE(n,H)\log^2 n)$ and hence
$\pc(n,H)=O(\pE(n,H)\log^3 n)$, using a rooted-extension
decomposition. They also established the conjectured single-logarithm
bound when $\pE(n,H)<1/(3n)$.
In further work~\cite[Theorems~1.4--1.5]{DKP25Cases}, they proved
$N(H,F)\le\EE N(G(n,\min\{1,Lq\}),F)$ whenever all subgraphs of $H$ have
expected count at least one at density $q$, for test graphs $F$ that
are cliques, cycles, or bounded-degree trees. The constant $L$ is
universal for cliques and cycles and depends on the degree bound
for trees.
Tran~\cite{Tran26} then obtained
$q_f(\mathcal F_H)=O(\pE(n,H)\log(2e(H)))$, reducing the general
containment loss to $O(\log^2(2e(H)))$. He removed the loss in the
comparison of expectation thresholds for trees and for graphs whose
average degree is at least the logarithm of their maximum degree,
thereby proving the second conjecture for these classes.
His ratio-maximizing cores give the conditional graph-extension
distributions used here.

Tran~\cite[Section~4]{Tran26} also proposed a coupling conjecture
that would absorb successive conditional extensions while paying the
logarithmic sampling cost once. We prove the required containment
bound through a tree reduction: we keep all possible extension histories
in a tree and reduce every level with the same random set. Disjointness along a
history supplies the precise control of dependence that this reuse
requires.

\subsection{Proof strategy and method ancestry}

Start with a nonempty edge subset $I$ of $H$. Among its subgraphs,
choose $S$ maximizing the probability that a fixed copy $S_0$ on
$[n]$ lies in a uniform copy of $I$, after division by a suitable
exponential weight in $e(S)$.
The first-moment constraints force $S$ to be smaller than $I$ by
a fixed factor. Maximality also bounds the probability that a prescribed
set $J$ of added edges appears, conditional on any copy of $S$, by
$(C_0\pE(n,H))^{|J|}$ for a universal constant $C_0$.
Repeating the choice gives nested subgraphs.
The tree branches over all their possible copies; its arc labels
are the added edges. Those labels are disjoint along each path,
and their sizes increase geometrically. This maximization argument
is the graph-specific spread-link construction of Tran~\cite[Section~3]{Tran26}.
Its conditioning step has antecedents in the large-link argument of
Lovett, Solomon, and Zhang~\cite[Claim~29]{LSZ19} and the spread-link
extraction from a maximal violating core in
Alweiss, Lovett, Wu, and Zhang~\cite[Lemma~2.4]{ALWZ21}.

For one distribution of labels, we use the Bayesian resampling
construction of Mossel, Niles-Weed, Sun, and Zadik~\cite{MNSZ25}.
Plant a source label into an independent random set $W$, sample a
target label from the posterior distribution given their union, and
retain the part of the target outside $W$. This remainder is also a
subset of the source. Section~\ref{sec:resampling} proves that large
remainders are exponentially unlikely and bounds the change of
measure on coordinates outside the target.

Section~\ref{sec:trees} applies this transition simultaneously at
every node of the fixed tree. A recursively defined potential bounds
the probability that too many continuations fail. Below a target
arc, all labels avoid that arc's label, so the local change-of-measure
estimate applies to the entire descendant potential. This is the
step that permits one $W$ to halve every level's capacity while
preserving spread at a controlled conditioning cost. Iteration needs
only logarithmically many successful draws. Finally,
Section~\ref{sec:graphs} constructs the graph hierarchy, with the
ordinary-copy normalization explicit, and applies the tree covering
theorem. All these arguments are proved below.

Throughout, $\log$ without a subscript is the natural logarithm.
For a finite set $X$ and $0\le\rho\le1$, let $\mu_\rho$ be the law
of a random subset of $X$ that includes each element independently
with probability $\rho$.
